%! Sine and Cosine Function Values — Unit 3, Lesson 3.3 — Guided Notes
\documentclass[10pt]{article}
\usepackage{precalculus-article}
\usepackage{precalculus-boxes}
\newcommand{\TallMath}[1]{$\displaystyle #1\rule[-1.4em]{0pt}{3.2em}$}

\begin{document}

\pageheader{Unit 3, Lesson 3.3}{Guided Notes}
\namedateperiod

\vspace{0.1in}

\begin{objectivebox}
By the end of this lesson, I will be able to\ldots
\begin{enumerate}[itemsep=3pt]
  \item \textbf{LO 3.3.A:} Determine the coordinates of a point on a circle centered at the origin using $P = (r\cos\theta,\, r\sin\theta)$.
  \item Use 45-45-90 and 30-60-90 triangle geometry to find exact sine and cosine values at multiples of $\pi/6$ and $\pi/4$.
  \item Apply quadrant sign conventions to find correct exact values for any standard-position angle.
\end{enumerate}
\end{objectivebox}

\vspace{0.08in}

\begin{vocabbox}
\termblanklong{reference angle}
\termblanklong{exact value}
\termblanklong{45-45-90 triangle}
\termblanklong{30-60-90 triangle}
\end{vocabbox}

\vspace{0.08in}

\begin{hookbox}
\textbf{GPS Satellite Problem:} A GPS satellite orbits at a radius of 26{,}560~km from Earth's center.
We model the orbit as a circle centered at the origin.
After the satellite rotates $\theta = \pi/3$ radians from its starting position, what are its exact $(x, y)$ coordinates?

\smallskip
What do you think we need to know to answer this?

\vspace{0.4in}
\end{hookbox}

\vspace{0.08in}

%-----------------------------------------------------------------------
\begin{notesbox}{Part A --- Exact Values from Special Triangles (EK 3.3.A.2)}

\textbf{45-45-90 triangle} (hypotenuse $= 1$):
\[
  \text{Each leg} = \underline{\hspace{1.5cm}} \qquad
  \cos\!\left(\frac{\pi}{4}\right) = \underline{\hspace{1.5cm}} \qquad
  \sin\!\left(\frac{\pi}{4}\right) = \underline{\hspace{1.5cm}}
\]

\vspace{0.1in}

\textbf{30-60-90 triangle} (hypotenuse $= 1$):
\[
  \text{Short leg} = \underline{\hspace{1.2cm}} \quad
  \text{Long leg} = \underline{\hspace{1.2cm}} \quad
  \cos\!\left(\frac{\pi}{6}\right) = \underline{\hspace{1.2cm}} \quad
  \sin\!\left(\frac{\pi}{6}\right) = \underline{\hspace{1.2cm}}
\]
\[
  \cos\!\left(\frac{\pi}{3}\right) = \underline{\hspace{1.5cm}} \qquad
  \sin\!\left(\frac{\pi}{3}\right) = \underline{\hspace{1.5cm}}
\]

\vspace{0.12in}

\textbf{Summary table} --- fill in the exact values:

\renewcommand{\arraystretch}{2.0}
\begin{tabularx}{\linewidth}{@{} >{\bfseries}l *{5}{>{\centering\arraybackslash}X} @{}}
\toprule
$\theta$ (rad) & $0$ & $\dfrac{\pi}{6}$ & $\dfrac{\pi}{4}$ & $\dfrac{\pi}{3}$ & $\dfrac{\pi}{2}$ \\
\midrule
$\sin\theta$ & \underline{\hspace{0.9cm}} & \underline{\hspace{0.9cm}} & \underline{\hspace{0.9cm}} & \underline{\hspace{0.9cm}} & \underline{\hspace{0.9cm}} \\
$\cos\theta$ & \underline{\hspace{0.9cm}} & \underline{\hspace{0.9cm}} & \underline{\hspace{0.9cm}} & \underline{\hspace{0.9cm}} & \underline{\hspace{0.9cm}} \\
\bottomrule
\end{tabularx}
\renewcommand{\arraystretch}{1}

\end{notesbox}

\vspace{0.08in}

%-----------------------------------------------------------------------
\begin{notesbox}{Part B --- Points on a Circle of Radius $r$ (EK 3.3.A.1)}

If a point $P$ lies on a circle of radius $r$ centered at the origin, at angle $\theta$ in standard position, then:
\[
  P = \bigl(\underline{\hspace{2cm}},\; \underline{\hspace{2cm}}\bigr)
\]

\vspace{0.08in}
\textbf{GPS Example:} $r = 26{,}560$~km, $\theta = \pi/3$

\[
  x = 26{,}560 \cdot \cos\!\left(\frac{\pi}{3}\right) = 26{,}560 \cdot \underline{\hspace{1.2cm}} = \underline{\hspace{2.4cm}} \text{ km}
\]
\[
  y = 26{,}560 \cdot \sin\!\left(\frac{\pi}{3}\right) = 26{,}560 \cdot \underline{\hspace{1.2cm}} = \underline{\hspace{2.4cm}} \text{ km}
\]

\end{notesbox}

\vspace{0.08in}

%-----------------------------------------------------------------------
\begin{notesbox}{Part C --- Signs in Each Quadrant}

\renewcommand{\arraystretch}{1.8}
\begin{tabularx}{\linewidth}{@{} >{\centering\arraybackslash}X | >{\centering\arraybackslash}X @{}}
\textbf{Quadrant II} ($x < 0,\; y > 0$) & \textbf{Quadrant I} ($x > 0,\; y > 0$) \\
$\cos\theta$ \underline{\hspace{0.9cm}}, $\sin\theta$ \underline{\hspace{0.9cm}} &
$\cos\theta$ \underline{\hspace{0.9cm}}, $\sin\theta$ \underline{\hspace{0.9cm}} \\
\hline
\textbf{Quadrant III} ($x < 0,\; y < 0$) & \textbf{Quadrant IV} ($x > 0,\; y < 0$) \\
$\cos\theta$ \underline{\hspace{0.9cm}}, $\sin\theta$ \underline{\hspace{0.9cm}} &
$\cos\theta$ \underline{\hspace{0.9cm}}, $\sin\theta$ \underline{\hspace{0.9cm}} \\
\end{tabularx}
\renewcommand{\arraystretch}{1}

\smallskip
Mnemonic: \textbf{A}ll \textbf{S}tudents \textbf{T}ake \textbf{C}alculus
(QI: All positive; QII: Sine positive; QIII: Tangent positive; QIV: Cosine positive).

\end{notesbox}

\vspace{0.08in}

%-----------------------------------------------------------------------
\begin{practicebox}
\textbf{Guided Practice 1.} Find the exact coordinates of the point $P$ on a circle of radius 3 at angle $\theta = \pi/3$.

\vspace{0.6in}

$P = \bigl($\underline{\hspace{1.5cm}}$,$\underline{\hspace{1.5cm}}$\bigr)$

\vspace{0.16in}

\textbf{Guided Practice 2.} A point $P = \bigl(-\tfrac{\sqrt{3}}{2},\, \tfrac{1}{2}\bigr)$ lies on the unit circle.
Identify the angle $\theta \in [0, 2\pi)$ in standard position. Justify using the sign of each coordinate.

\vspace{0.7in}

$\theta = $ \underline{\hspace{2cm}}
\end{practicebox}

\end{document}
